\subsection{THE LARGEST NILPOTENT IDEAL OF A LIE ALGEBRA}

Let \(\mathfrak{g}\) be a Lie algebra and \(\mathfrak{a}\) an ideal of \(\mathfrak{g}\) . For \(\mathfrak{a}\) to be nilpotent, it is necessary and sufficient that, for all \(x \in \mathfrak{a}\) , \(\mathrm{ad}_{\mathfrak{g}}x\) be nilpotent; the condition is obviously sufficient and is necessary, for, if \(\mathfrak{a}\) is nilpotent and \(x \in \mathfrak{a}\) , \(\mathrm{ad}_{\mathfrak{a}}x\) is nilpotent and \(\mathrm{ad}_{\mathfrak{a}}x\) maps \(\mathfrak{g}\) into \(\mathfrak{a}\) , hence \(\mathrm{ad}_{\mathfrak{g}}x\) is nilpotent. Then Proposition 4 applied to the adjoint representation of \(\mathfrak{g}\) gives the following result:

PROPOSITION 6. Let \(\mathfrak{g}\) be a Lie algebra and E the associative subalgebra of \(\mathcal{L}(\mathfrak{g})\) generated by 1 and the \(\operatorname{ad}_{\mathfrak{g}}x(x\in \mathfrak{g})\) . Let R be the Jacobson radical of E.

\begin{enumerate}
\def\labelenumi{(\alph{enumi})}
\item
  The set \(\mathfrak{n}\) of \(y\in \mathfrak{g}\) such that \(\operatorname {ad}_{\mathfrak{g}}y\in \mathbb{R}\) is the largest nilpotent ideal of \(\mathfrak{g}\) .
\item
  It is orthogonal to g under the Killing form.
\end{enumerate}

It should be noted that g/n can have non-zero nilpotent ideals.

